\section{Extrema and measurable ellipticity}
\label{uni-sec:anchors}

Let $X$ be a connected compact complex manifold of dimension $n\geq1$, let $\omega$ be a Hermitian form on $X$, and fix a smooth positive volume form $\mathrm{d}V$. For real potentials $u,v$ of the regularity specified below, put $A=\omega+\ddc u$, $B=\omega+\ddc v$, $w=u-v$ and $T=\sum_{j=0}^{n-1}A^j\wedge B^{n-1-j}$. Products of measurable coefficients are interpreted almost everywhere. Suppose $u,v\in C^{1,1}(X,\mathbb R)$, $A,B\geq0$ almost everywhere, and
\begin{equation}
  A^n=B^n=f\,\mathrm{d}V,\qquad f\in C^0(X),\qquad f\geq0,
  \label{uni-eq:density}
\end{equation}
where the coefficient products are understood almost everywhere.  Fix a
smooth positive $(1,1)$-form $\beta$, put $H=A+B$, and set
\begin{equation}
 E_{\rm unif}
 =\left\{x:\begin{array}{l}
 \text{there are a neighborhood }V\ni x\text{ and }c>0\\
 \text{with }H\geq c\beta\text{ almost everywhere on }V
 \end{array}\right\}.
 \label{uni-eq:Eunif}
\end{equation}
This open set is independent of $\beta$ and of the representatives of the
measurable coefficients.

\begin{theorem}\label{uni-thm:anchors}
Let $u,v\in C^{1,1}(X)$, assume that $A,B\geq0$ almost everywhere, and
assume \eqref{uni-eq:density}.  If $\chi\in C^2(X,\mathbb R)$ satisfies
$C:=\omega+\ddc\chi>0$, put
\[
  P_\chi=\{x\in X:f(x)\mathrm{d}V_x\geq C_x^n\},\qquad U=\{f>0\}.
\]
Then $P_\chi$ is a nonempty compact subset of $U$, and
$U\subset E_{\rm unif}$.  There are components $G_+,G_-$ of
$E_{\rm unif}$, each meeting $P_\chi$, such that
\begin{equation}
 w\equiv\max_Xw\quad\hbox{on }\overline{G_+},\qquad
 w\equiv\min_Xw\quad\hbox{on }\overline{G_-}.
 \label{uni-eq:extrema-components}
\end{equation}
If $w$ is nonconstant, these closures are disjoint.  Consequently $w$ is
constant if any one of the following holds:
\begin{enumerate}
\item every two components of $E_{\rm unif}$ meeting $P_\chi$ have
      intersecting closures;
\item $E_{\rm unif}$ is connected;
\item $E_{\rm unif}$ is dense and has at most two components.
\end{enumerate}
The same three conclusions hold with $U$ in place of $E_{\rm unif}$.
In particular, uniqueness holds when $\{f=0\}$ is a nonempty connected
embedded $C^1$ real hypersurface without boundary.  No orientability or
coorientability assumption is needed.
\end{theorem}

We begin with two local lemmas.

\begin{lemma}\label{uni-lem:goodcontact}
In Euclidean space $\mathbb R^m$, let $B(p,r)$ be the ball of center $p$ and radius $r>0$, and let $I$ be the identity matrix. Let $g\in C^{1,1}$ near $\overline{B(p,r)}\subset\mathbb R^m$, suppose
$g\leq g(p)$ on the ball, and let $E$ be any full-measure subset on
which $g$ has a second-order Taylor expansion.  For every $\delta>0$ and
$0<\rho<r$, there is $x\in E\cap B(p,\rho)$ such that
\[
    D^2g(x)\leq2\delta I.
\]
\end{lemma}

\begin{proof}
Put $g_\delta(x)=g(x)-\delta|x-p|^2$.  The gradient map
$F=\nabla g_\delta$ is Lipschitz on the ball, so it maps the null set
$N=\overline{B(p,r)}\setminus E$ to a null set.  Choose
$s\in B(0,\delta\rho)\setminus F(N)$ and maximize
$g_\delta(x)-s\cdot(x-p)$ on the closed ball.  At any point with
$|x-p|\geq\rho$, its value is strictly smaller than its value at $p$,
because
\[
 g(x)-g(p)-\delta|x-p|^2-s\cdot(x-p)
 \leq-\delta|x-p|^2+|s|\,|x-p|<0.
\]
Thus a maximizer $x$ lies in $B(p,\rho)$.  Its first derivative gives
$F(x)=s$, so $x\notin N$.  The second-order Taylor expansion at this
upper contact point gives $D^2g(x)\leq2\delta I$.
The null-set assertion for a locally Lipschitz gradient follows by a
countable cover by sets on which it is Lipschitz.  The full-measure set of
second-order Taylor points exists since the first derivatives of a
$C^{1,1}$ function are Lipschitz.
\end{proof}

\begin{lemma}\label{uni-lem:measSMP}
Let $a=(a^{ij})$ be a measurable symmetric matrix field on a connected
domain $\Omega\subset\mathbb R^m$. For each compact $K\Subset\Omega$, suppose there are constants $0<\lambda_K\leq\Lambda_K<\infty$ such that
\[
   \lambda_K I\leq a(x)\leq\Lambda_KI
\]
almost everywhere on $K$, where $I$ is the identity matrix. Repeated indices $i,j=1,\ldots,m$ are summed, and $z_{ij}=\partial_i\partial_jz$. If
$z\in W^{2,\infty}_{\rm loc}(\Omega)$,
$a^{ij}z_{ij}\geq0$ almost everywhere, and $z$ attains its supremum at an
interior point, then $z$ is constant.
\end{lemma}

\begin{proof}
We first prove strict weak comparison on a relatively compact domain:
if $b\in C^{1,1}$, $Lb:=a^{ij}b_{ij}\geq\kappa>0$ almost everywhere,
and $b\leq M$ on the boundary, then $b\leq M$ inside.  Otherwise an
interior maximum lies in a closed coordinate ball on which
$a\leq\Lambda I$.  Apply Lemma~\ref{uni-lem:goodcontact} to the full-measure
set where both the differential inequality and the Taylor expansion hold.
Choosing $2\delta m\Lambda<\kappa$ gives the contradiction
$\kappa\leq Lb\leq2\delta\operatorname{tr}a\leq2\delta m\Lambda$.

Take the $C^{1,1}$ representative of $z$, write $M=\sup_\Omega z$,
and suppose $K=\{z=M\}$ is nonempty and proper.  Near a relative boundary
point of $K$, choose $y\notin K$ so that the ball of radius
$r=\operatorname{dist}(y,K)>0$ is compactly contained in $\Omega$ and
touches $K$ at $x_0$.  On the closed annulus
$r/2\leq|x-y|\leq r$, use fixed bounds $\lambda I\leq a\leq\Lambda I$
and set
\[
 b(x)=\mathrm{e}^{-\alpha|x-y|^2}-\mathrm{e}^{-\alpha r^2}.
\]
For sufficiently large $\alpha$,
\[
 Lb=\mathrm{e}^{-\alpha|x-y|^2}
       \bigl(4\alpha^2 a^{ij}(x_i-y_i)(x_j-y_j)
                       -2\alpha\operatorname{tr}a\bigr)\geq\kappa>0
\]
almost everywhere on the annulus.  On its inner sphere $z<M$ uniformly;
on its outer sphere $b=0$.  Choose $\eta>0$ small enough that
$z+\eta b\leq M$ on both spheres.  Strict weak comparison gives this
inequality throughout the annulus.  At $x_0$, however, $\nabla z=0$
and $b$ increases to first order in the inward radial direction, which
contradicts that inequality.  Thus $K$ has no relative boundary and
connectedness proves the assertion.  This argument uses no derivative of
the measurable coefficients.
\end{proof}

At a point where $A,B$ are semipositive, define the cofactor quadratic form on real cotangent vectors
\begin{equation}
 Q_x(\xi)\mathrm{d}V_x
   =\mathrm{i}\,\xi^{1,0}\wedge\overline{\xi^{1,0}}\wedge T_x .
 \label{uni-eq:Q}
\end{equation}
Mixed products of semipositive $(1,1)$-forms are semipositive in this
covector-testing order.  Expanding $H^{n-1}$, $H=A+B$, gives
\begin{equation}
 2^{-(n-1)}H^{n-1}\leq T\leq H^{n-1}.
 \label{uni-eq:cofactor-comparison}
\end{equation}
At a point, choose a $\beta$-unitary coframe that diagonalizes $H$, and write $\lambda_1,\ldots,\lambda_n$ for its eigenvalues. In that coframe the
coefficient of
$\mathrm{i}\alpha\wedge\bar\alpha\wedge H^{n-1}$, for
$\alpha=\sum_{k=1}^n b_k\,\mathrm{d}z_k$, with complex coefficients $b_k$, is a positive dimensional constant times
\begin{equation}
  \sum_{k=1}^n |b_k|^2\prod_{\ell\ne k}\lambda_\ell.
 \label{uni-eq:cofactor-eigen}
\end{equation}
For $n\geq2$, \eqref{uni-eq:cofactor-comparison}--\eqref{uni-eq:cofactor-eigen}
show that $Q$ is locally uniformly elliptic exactly where $H$ has a
local essential positive lower bound.  The reverse implication uses the
local upper bound for $H$: testing each eigendirection bounds every
$\prod_{\ell\ne k}\lambda_\ell$ below, and choosing $k$ different from
an index of a smallest eigenvalue bounds that eigenvalue below.  For $n=1$
one has $T=1$, so the linearized equation is elliptic everywhere.

\begin{proof}[of Theorem~\ref{uni-thm:anchors}]
We first locate the extrema.  Fix $s>0$ and maximize
\begin{equation}
  w-s(v-\chi).
 \label{uni-eq:pertmax}
\end{equation}
For smooth data, the complex Hessian inequality at a maximizer is
\begin{equation}
  (1+s)B\geq A+sC.
 \label{uni-eq:matrixcontact}
\end{equation}
The positive term $sC$ first gives $B>0$, and determinant equality then
gives $A>0$.  In a $B$-unitary frame put
$J=B^{-1/2}AB^{-1/2}$, $K=B^{-1/2}CB^{-1/2}$.  Since $\det J=1$,
arithmetic--geometric mean and \eqref{uni-eq:matrixcontact} give
\[
 n+s\operatorname{tr}K\leq\operatorname{tr}J+s\operatorname{tr}K
 \leq n(1+s).
\]
Consequently $\operatorname{tr}K\leq n$, $B\geq C/n$, and
$\det K\leq(\operatorname{tr}K/n)^n\leq1$.  In particular $B^n\geq C^n$.

For $C^{1,1}$ potentials, fix a maximizer $p$ of
\eqref{uni-eq:pertmax} and apply Lemma~\ref{uni-lem:goodcontact} on a coordinate
ball, choosing radii tending to zero and errors $\delta_j\downarrow0$.
Include in its full-measure set the points where both Hessians, their
semipositivity, and their common determinant are defined.  With
$\beta_0=\ddc|z|^2$, the resulting points $x_j\to p$ satisfy
\[
 A(x_j)+s C_j(x_j)\leq(1+s)B(x_j),\qquad
 C_j=C-(\delta_j/s)\beta_0>0
\]
for large $j$.  The matrix argument just given yields
$f(x_j)\mathrm{d}V_{x_j}\geq C_j(x_j)^n$.  Continuity of $f,\mathrm{d}V,C$, not
continuity of the Hessians, gives $p\in P_\chi$.  Thus every maximizer of
the perturbed function lies in $P_\chi$.  Letting $s\downarrow0$
places an actual maximizer of $w$ in $P_\chi$.  Applying the same argument
to $-w$ places a minimizer there as well.  Compactness and positivity of
$C$ show that $P_\chi$ is compact, nonempty, and contained in $U$.

On a relatively compact ball in $U$, the determinants of $A$ and $B$
have a positive lower bound, while their coefficients have an essential
upper bound.  Every eigenvalue is therefore bounded below by the determinant
lower bound divided by a fixed power of the upper bound.  Thus
$U\subset E_{\rm unif}$.

By \eqref{uni-eq:linearization}, $w$ solves a measurable nondivergence equation
whose principal quadratic form is a fixed positive multiple of $Q$.  On
each relatively compact coordinate ball in $E_{\rm unif}$, its coefficients
are measurable and uniformly elliptic.  Lemma~\ref{uni-lem:measSMP} makes $w$
constant on the component containing its maximum found above, and makes $-w$
constant on the component containing its minimum found above.  Continuity
extends the two constant values to the closures, proving
\eqref{uni-eq:extrema-components}.  Distinct extrema force disjoint closures.

The same argument applies on $U$ itself: it contains the extremal points in $P_\chi$, and the cofactor is locally uniformly elliptic there. This proves the assertions with $U$ in place of $E_{\rm unif}$.

The component criteria are immediate except for the dense two-component
case.  If an open dense subset of connected $X$ has at most two components
with disjoint closures, those closures are disjoint nonempty closed sets
whose union is $X$, contradicting connectedness.  Finally, when
$Z=\{f=0\}$ is a connected embedded real hypersurface, flattening charts
show that all local sides belong to one fixed unordered pair of components
of $X\setminus Z$.  Hence $U$ has at most two components; it is dense
because a hypersurface has empty interior.  This proves the last assertion.
\end{proof}

